\documentclass[11pt,a4paper]{article}

\usepackage[a4paper,margin=2.5cm]{geometry}
\usepackage[utf8]{inputenc}
\usepackage[T1]{fontenc}
\usepackage{lmodern}
\usepackage{titlesec}
\usepackage{fancyhdr}
\usepackage{enumitem}

\titleformat{\section}{\normalfont\Large\bfseries}{\thesection}{0.6em}{}
\titlespacing*{\section}{0pt}{1.1em}{0.5em}

\pagestyle{fancy}
\fancyhf{}
\renewcommand{\headrulewidth}{0pt}
\fancyfoot[C]{\small YANFLÈ --- Orange Summer Challenge 2026 --- \thepage}

\setlength{\parindent}{0pt}
\setlength{\parskip}{0.45em}

\begin{document}

\begin{center}
  {\Huge\bfseries YANFLÈ}\\[6pt]
  {\large En dioula : ``regarde ici''}\\[4pt]
  {\itshape Tech Impact --- Orange Summer Challenge 2026}
\end{center}

\vspace{1.5em}

Le projet consiste à donner aux caméras de surveillance la capacité de comprendre ce qu'elles filment, et non plus seulement de détecter un mouvement. Au lieu du détecteur de mouvement traditionnel, qui signale une agitation sans savoir ce qu'elle signifie, le système observe en temps réel et reconnaît des situations : une présence prolongée et suspecte, un comportement qui précède un vol, un cambriolage ou une agression. Il donne l'alerte avant que l'acte ne se produise, sur la base d'une véritable prédiction du danger.

\section{Le problème}
Les caméras de surveillance actuelles sont passives. Elles enregistrent, mais elles n'analysent rien de ce qui se passe réellement devant elles. Le détecteur de mouvement classique ne fait pas la différence entre un événement anodin et une situation dangereuse : il réagit à un déplacement, jamais à un comportement.

Cette passivité a une conséquence directe : la caméra ne sert qu'à constater les faits après coup, une fois le crime, le cambriolage ou l'agression déjà commis. Elle ne prévient personne, elle ne fait que filmer.

\section{La solution}
L'idée est de remplacer cette surveillance passive par un système qui regarde à la place des gens et qui comprend ce qu'il observe, en s'appuyant sur les caméras déjà installées.

Concrètement, l'approche repose sur quatre principes :

\begin{enumerate}[leftmargin=1.4em, itemsep=2pt]
  \item \textbf{Observer le comportement, pas l'identité.} Le système s'intéresse à la posture, aux déplacements et aux gestes des personnes filmées, pas à la reconnaissance de leur visage.
  \item \textbf{Reconnaître les signaux qui précèdent un acte.} Le modèle est entraîné à repérer les comportements caractéristiques d'une situation à risque avant qu'elle ne dégénère, plutôt que de simplement constater qu'il y a du mouvement.
  \item \textbf{Alerter un humain, ne jamais décider seul.} Le système ne fait que signaler une situation suspecte à un opérateur, qui reste seul décisionnaire de la suite à donner.
  \item \textbf{Adapter la réponse à la gravité de la situation.} Une simple notification suffit pour un signal faible ; pour une situation critique, le système va jusqu'à déclencher un appel automatisé vers la personne responsable, afin de s'assurer que l'alerte a bien été prise en compte.
\end{enumerate}

L'ambition du projet est de faire passer la vidéosurveillance d'un rôle de témoin, qui ne sert qu'à prouver qu'un acte a eu lieu, à un rôle de vigie, capable d'aider à l'empêcher.

\vspace{1em}
\noindent\textit{À ce stade, il s'agit d'une idée de projet à construire : la conception, l'entraînement du modèle et les tests restent entièrement à faire.}

\end{document}
